% !TeX program = pdflatex
% corpus/docs/analise_exercicios.tex — ticks AR1–AR25 (supremo → cinco vias).
%
% Prosa canónica; testemunhas em analise_resolve() (banco/conversa.c).
% Espinha: HIPÓTESES → DEFINIÇÃO → TRANSIÇÃO → LEI → TESTEMUNHA → CONCLUSÃO → VOLTA.
%
\documentclass[11pt,a4paper]{article}
\usepackage[utf8]{inputenc}
\usepackage[T1]{fontenc}
\usepackage[portuguese]{babel}
\usepackage{amsmath,amssymb,amsthm}
\usepackage[margin=2.4cm]{geometry}
\usepackage{booktabs}
\usepackage{xcolor}
\usepackage[hidelinks]{hyperref}

\newcommand{\code}[1]{\texttt{\small #1}}
\newcommand{\medido}[1]{\par\medskip\noindent{\small\color{gray}\textbf{Medido.} #1}\par\medskip}

\begin{document}
\title{Análise real --- exercícios\\
\large AR1--AR25: supremo $\to$ cinco vias}
\author{Tiffany --- corpus vivo}
\date{}
\maketitle

\begin{abstract}
\noindent Vinte e cinco exercícios (\code{AR25[]} em \code{banco/conversa.c}) ---
este ficheiro fecha o catálogo. Operação ao vivo: nomes do catálogo via
\code{responde} / \code{analise N}.
\end{abstract}

\section{como se le este documento}
Espinha dos sete slots. Lotes: AR1--AR8 (supremo\ldots Bolzano);
AR9--AR16 ($\varepsilon$-$N$\ldots Weierstrass);
AR17--AR25 (derivada\ldots cinco vias).
Títulos que colidem com Cálculo I levam sufixo distintivo
(\code{derivada limite e quociente}, \code{valor medio na analise},
\code{rolle e o gume}).

%======================================================================
\section{supremo}\label{sec:supremo}
\textbf{Enunciado.} $\sup$ e $\inf$: a EXISTÊNCIA separada do atingimento.

\begin{enumerate}\itemsep2pt
\item \textbf{DEFINIÇÃO.} $A\subseteq\mathbb{Q}$ não vazio e limitado; $\varepsilon>0$
  racional.
\item \textbf{TRADUÇÃO.} $\sup A$ existe $\not\Rightarrow\sup A\in A$.
\item \textbf{OPERAÇÃO.} Lista finita confunde os três estados; $A=\{q>0:q^2<2\}$
  separa-os.
\item \textbf{LEI.} Ordem em $\mathbb{Q}$ por produto cruzado --- dois inteiros.
\item \textbf{TESTEMUNHA.} Existir / ser atingido / ser racional --- três estados
  distintos no mesmo $A$.
\item \textbf{CONCLUSÃO.} Existência do supremo é propriedade do ESPAÇO --- e
  $\mathbb{Q}$ não a tem.
\item \textbf{VOLTA.} O que $\mathbb{R}$ acrescenta é exactamente o supremo que falta
  (cinco formas na fala~25).
\end{enumerate}
\medido{\code{analise\_resolve(1)}: lista finita atinge; $A$ de $\sqrt{2}$ não.}

%----------------------------------------------------------------------
\section{arquimediana}\label{sec:arquimediana}
\textbf{Enunciado.} Dado $\varepsilon$, achar $n$ com $1/n<\varepsilon$ ---
construtivo, não «existe».

\begin{enumerate}\itemsep2pt
\item \textbf{DEFINIÇÃO.} $\varepsilon=p/q>0$ racional.
\item \textbf{TRADUÇÃO.} $\forall\varepsilon>0\ \exists n:\ 1/n<\varepsilon$.
\item \textbf{OPERAÇÃO.} $n=\lfloor q/p\rfloor+1$ --- ESTE $n$, verificado.
\item \textbf{LEI.} Ordem em $\mathbb{Q}$ (produto cruzado).
\item \textbf{TESTEMUNHA.} $n$ exibido para cada $\varepsilon$ da malha.
\item \textbf{CONCLUSÃO.} $\mathbb{Q}$ é arquimediano --- nenhum racional
  «infinitamente pequeno».
\item \textbf{VOLTA.} Existência construtiva $>$ existência só provada.
\end{enumerate}
\medido{\code{analise\_resolve(2)} / \code{an\_arquimedes}: $n$ para cada
$\varepsilon$; $0$ falhas.}

%----------------------------------------------------------------------
\section{densidade}\label{sec:densidade-q}
\textbf{Enunciado.} Um racional entre quaisquer dois: o ponto médio.

\begin{enumerate}\itemsep2pt
\item \textbf{DEFINIÇÃO.} $a<b$ racionais.
\item \textbf{TRADUÇÃO.} $a<(a+b)/2<b$.
\item \textbf{OPERAÇÃO.} Ponto médio exacto em $\mathbb{Q}$ --- uma linha, toda a
  densidade.
\item \textbf{LEI.} Fecho de $\mathbb{Q}$ para a média.
\item \textbf{TESTEMUNHA.} Médio exibido em $[1/3,1/2]$ e em $[0,1/1000]$.
\item \textbf{CONCLUSÃO.} $\mathbb{Q}$ denso em si próprio.
\item \textbf{VOLTA.} Entre quaisquer dois há racionais --- e também irracionais
  (fala seguinte).
\end{enumerate}
\medido{\code{analise\_resolve(3)} / \code{an\_entre}: médio sempre cabe.}

%----------------------------------------------------------------------
\section{irracionais densos}\label{sec:irracionais-densos}
\textbf{Enunciado.} $q+\sqrt{2}/n$ dentro de qualquer intervalo --- e o que NÃO se
avalia.

\begin{enumerate}\itemsep2pt
\item \textbf{DEFINIÇÃO.} $a<b$ racionais; candidato $q+\sqrt{2}/n$.
\item \textbf{TRADUÇÃO.} Irracional no intervalo, sem avaliar $\sqrt{2}$.
\item \textbf{OPERAÇÃO.} Desigualdade AO QUADRADO: $(b-q)^2 n^2>2$ --- exacta em
  $\mathbb{Q}$.
\item \textbf{LEI.} Paridade / não-raiz em $\mathbb{Q}$ --- não comparação decimal.
\item \textbf{TESTEMUNHA.} Quadrado comparado; $\sqrt{2}$ fica por tirar.
\item \textbf{CONCLUSÃO.} Irracionais densos; entre dois racionais há dos dois tipos.
\item \textbf{VOLTA.} Escrever «$\sqrt{2}\approx1{,}414$» traz o decimal --- e o
  decimal não é o objecto.
\end{enumerate}
\medido{\code{analise\_resolve(4)}: desigualdade ao quadrado; sem avaliar a raiz.}

%----------------------------------------------------------------------
\section{dedekind}\label{sec:dedekind}
\textbf{Enunciado.} O corte de $\sqrt{2}$, e as DUAS propriedades verificadas.

\begin{enumerate}\itemsep2pt
\item \textbf{DEFINIÇÃO.} Sucessão racional que aproxima $\sqrt{2}$.
\item \textbf{TRADUÇÃO.} $A=\{q\le0\}\cup\{q:q^2<2\}$, $B=$ o resto.
\item \textbf{OPERAÇÃO.} $A$ fechado para baixo; $A$ sem máximo --- dado $a\in A$,
  exibe-se $a'>a$ em $A$ ($a'=(2a+2)/(a+2)$).
\item \textbf{LEI.} Nenhum racional tem quadrado $2$ (paridade dos expoentes).
\item \textbf{TESTEMUNHA.} PASSO verificado em muitos $a\in A$ independentes --- não
  só a órbita.
\item \textbf{CONCLUSÃO.} $(A,B)$ corte legítimo; fronteira não racional.
\item \textbf{VOLTA.} Completamento: $\mathbb{R}$ é onde Cauchy e «ter limite»
  coincidem.
\end{enumerate}
\medido{\code{analise\_resolve(5)} / \code{an\_passo\_sobe}: $0$ falhas no passo.}

%----------------------------------------------------------------------
\section{cauchy em Q}\label{sec:cauchy-em-q}
\textbf{Enunciado.} De Cauchy sem usar o limite --- o gume, na assinatura.

\begin{enumerate}\itemsep2pt
\item \textbf{DEFINIÇÃO.} Mesma sucessão racional a aproximar $\sqrt{2}$.
\item \textbf{TRADUÇÃO.} $\forall\varepsilon>0\ \exists N:\ m,n>N\Rightarrow|x_m-x_n|<\varepsilon$.
\item \textbf{OPERAÇÃO.} Módulo compara $|x_m-x_n|$ e NUNCA menciona limite --- a
  sucessão aproxima-se DE SI.
\item \textbf{LEI.} Ser de Cauchy $\neq$ ter limite; $\mathbb{Q}$ separa as duas.
\item \textbf{TESTEMUNHA.} $N$ achado para cada $\varepsilon$; busca exaustiva do
  limite racional volta VAZIA.
\item \textbf{CONCLUSÃO.} Cauchy pertence à sucessão; limite depende do espaço.
\item \textbf{VOLTA.} Completude $=$ coincidência das duas afirmações.
\end{enumerate}
\medido{\code{analise\_resolve(6)} / \code{cy\_modulo}: Cauchy sim; limite em
$\mathbb{Q}$ não.}

%----------------------------------------------------------------------
\section{completude}\label{sec:completude}
\textbf{Enunciado.} Cauchy converge $\Leftrightarrow$ todo limitado tem supremo.

\begin{enumerate}\itemsep2pt
\item \textbf{DEFINIÇÃO.} Duas afirmações sobre $\mathbb{R}$.
\item \textbf{TRADUÇÃO.} Todo Cauchy converge $\Leftrightarrow$ todo limitado
  superiormente tem $\sup$.
\item \textbf{OPERAÇÃO.} Duas direcções À PARTE: ida (conjunto $\to$ sucessão);
  volta (sucessão $\to$ conjunto) --- mesma bisseção, dois sentidos.
\item \textbf{LEI.} Bisseção: motor das duas; muda o que se procura.
\item \textbf{TESTEMUNHA.} Encaixe medido; módulo de Cauchy para $\varepsilon$ fino.
\item \textbf{CONCLUSÃO.} As duas são a MESMA --- «completude» nomeia a coincidência.
\item \textbf{VOLTA.} Em $\mathbb{Q}$ falham JUNTAS: encaixe sem ponto, conjunto sem
  $\sup$.
\end{enumerate}
\medido{\code{analise\_resolve(7)}: $24$ bissecções + módulo; ambas fecham.}

%----------------------------------------------------------------------
\section{bolzano}\label{sec:bolzano}
\textbf{Enunciado.} Limitada tem subsucessão convergente --- procurada.

\begin{enumerate}\itemsep2pt
\item \textbf{DEFINIÇÃO.} $(x_n)$ limitada em $[1,2]$ --- convergentes de $\sqrt{2}$.
\item \textbf{TRADUÇÃO.} Limitada $\Rightarrow$ existe subsucessão convergente.
\item \textbf{OPERAÇÃO.} Bisseção: metade com mais termos; sempre há uma.
\item \textbf{LEI.} Encaixe: caixas encolhem com termos; em $\mathbb{R}$ o ponto é o
  limite.
\item \textbf{TESTEMUNHA.} Índices EXIBIDOS; caixa final; ponto não racional
  ($p^2=2q^2$ vazio).
\item \textbf{CONCLUSÃO.} Subsucessão sempre; limite só em $\mathbb{R}$ --- duas
  afirmações.
\item \textbf{VOLTA.} Bolzano--Weierstrass É forma da completude (fala~25).
\end{enumerate}
\medido{\code{analise\_resolve(8)} / \code{an\_subsuc}: índices e caixa; sem raiz
racional.}

%----------------------------------------------------------------------
\section{epsilon N}\label{sec:epsilon-n}
\textbf{Enunciado.} $(2n+1)/(n+3)\to2$ com o $N$ CONSTRUÍDO de $\varepsilon$.

\begin{enumerate}\itemsep2pt
\item \textbf{DEFINIÇÃO.} $x_n=(2n+1)/(n+3)$.
\item \textbf{TRADUÇÃO.} $|x_n-2|=5/(n+3)<\varepsilon$ para $n>5/\varepsilon-3$.
\item \textbf{OPERAÇÃO.} $N$ sai da CONTA em $\mathbb{Q}$ --- o limite aparece como o
  valor para o qual a diferença encolhe.
\item \textbf{LEI.} Completude de $\mathbb{R}$ (encaixe tem ponto).
\item \textbf{TESTEMUNHA.} $N$ construído e distância VERIFICADA para cada
  $\varepsilon$.
\item \textbf{CONCLUSÃO.} $x_n\to2$ com $N$ explícito.
\item \textbf{VOLTA.} Bolzano--Weierstrass: divergir ainda permite subsucessões
  convergentes.
\end{enumerate}
\medido{\code{analise\_resolve(9)}: $N=5/\varepsilon-3$; $0$ falhas.}

%----------------------------------------------------------------------
\section{duas subsucessoes}\label{sec:duas-subsucessoes}
\textbf{Enunciado.} Dois limites diferentes provam a divergência.

\begin{enumerate}\itemsep2pt
\item \textbf{DEFINIÇÃO.} $(x_n)$ limitada --- tipicamente $(-1)^n$.
\item \textbf{TRADUÇÃO.} Dois limites de subsucessão distintos $\Rightarrow$ diverge.
\item \textbf{OPERAÇÃO.} Pares $\to+1$, ímpares $\to-1$; critérios de cauda
  separam.
\item \textbf{LEI.} Completude: encaixe pode ter ponto; convergência da sucessão
  exige um só.
\item \textbf{TESTEMUNHA.} Par $(-1)^n$ vs $1/(n+1)$: limsup/liminf iguais só no
  segundo.
\item \textbf{CONCLUSÃO.} Divergência PROVADA por dois limites distintos.
\item \textbf{VOLTA.} Existência de subsucessão $\neq$ existência de limite da
  sucessão.
\end{enumerate}
\medido{\code{analise\_resolve(10)}: $(-1)^n$ nunca iguala caudas.}

%----------------------------------------------------------------------
\section{limsup}\label{sec:limsup}
\textbf{Enunciado.} Converge $\Leftrightarrow$ limsup $=$ liminf, e $(-1)^n$ é o
gume.

\begin{enumerate}\itemsep2pt
\item \textbf{DEFINIÇÃO.} $(x_n)$ limitada.
\item \textbf{TRADUÇÃO.} Converge $\Leftrightarrow\lim_n\sup_{k\ge n}x_k=\lim_n\inf_{k\ge n}x_k$.
\item \textbf{OPERAÇÃO.} Limsup/liminf $=$ $\sup$/$\inf$ da CAUDA --- existem sempre;
  convergência é a igualdade.
\item \textbf{LEI.} Caracterização completa, não só critério suficiente.
\item \textbf{TESTEMUNHA.} $(-1)^n$: $+1$ e $-1$ nunca se aproximam; $1/(n+1)$:
  distância encolhe.
\item \textbf{CONCLUSÃO.} Converge $\Leftrightarrow$ limsup $=$ liminf.
\item \textbf{VOLTA.} Mesmo divergindo, $(-1)^n$ tem subsucessões convergentes
  (Bolzano).
\end{enumerate}
\medido{\code{analise\_resolve(11)} / \code{an\_cauda}: gume vs convergente.}

%----------------------------------------------------------------------
\section{tres continuidades}\label{sec:tres-continuidades}
\textbf{Enunciado.} $\varepsilon$-$\delta$, sequencial e $f^{-1}(F)$ fechado --- as
três.

\begin{enumerate}\itemsep2pt
\item \textbf{DEFINIÇÃO.} $f(x)=x^2$, $a\in\mathbb{Q}$.
\item \textbf{TRADUÇÃO.} $\varepsilon$-$\delta$; sequencial; $f^{-1}(F)$ fechado.
\item \textbf{OPERAÇÃO.} Medem-se À PARTE: $\delta$ da factorização
  $|x^2-a^2|=|x-a|\,|x+a|$; sequencial com $x_n=a+1/n$; topológica sem distância.
\item \textbf{LEI.} $\delta=\min(1,\varepsilon/(2|a|+1))$.
\item \textbf{TESTEMUNHA.} Três colunas: $\delta$ serve; distâncias descem; forma
  topológica sem régua.
\item \textbf{CONCLUSÃO.} Coincidem em métricos; SÓ a terceira generaliza.
\item \textbf{VOLTA.} Irmão mais forte: continuidade uniforme (fala~13) --- ordem dos
  quantificadores.
\end{enumerate}
\medido{\code{analise\_resolve(12)} / \code{an\_delta}: três vias; $0$ falhas.}

%----------------------------------------------------------------------
\section{uniforme}\label{sec:uniforme-ar}
\textbf{Enunciado.} $1/x$ em $(0,1)$: contínua e NÃO uniformemente.

\begin{enumerate}\itemsep2pt
\item \textbf{DEFINIÇÃO.} $f(x)=1/x$ em $(0,1)$.
\item \textbf{TRADUÇÃO.} $\delta$ depende do PONTO, não só de $\varepsilon$.
\item \textbf{OPERAÇÃO.} Contínua: $\exists\delta$ por $x$; uniforme: $\exists\delta$
  para TODOS $x$ --- ordem dos quantificadores.
\item \textbf{LEI.} $\forall x\,\exists\delta$ contra $\exists\delta\,\forall x$.
\item \textbf{TESTEMUNHA.} De cada $\delta$: $x=\delta/2$, $y=\delta/3$; salto
  exacto $1/\delta$ --- nenhum $\delta$ serve para $\varepsilon=1$.
\item \textbf{CONCLUSÃO.} Contínua e NÃO uniformemente contínua.
\item \textbf{VOLTA.} Em $[a,b]$ compacto toda contínua é uniformemente contínua ---
  falta a $(0,1)$ a compacidade.
\end{enumerate}
\medido{\code{analise\_resolve(13)} / \code{an\_falha\_uniforme}: salto $1/\delta$.}

%----------------------------------------------------------------------
\section{compacto}\label{sec:compacto-ar}
\textbf{Enunciado.} Sucessão em compacto tem limite DENTRO dele.

\begin{enumerate}\itemsep2pt
\item \textbf{DEFINIÇÃO.} $K\subseteq\mathbb{R}$ compacto; $f$ contínua em $K$.
\item \textbf{TRADUÇÃO.} Toda sucessão em $K$ tem subsucessão convergente COM limite
  em $K$.
\item \textbf{OPERAÇÃO.} Limite fica DENTRO: compacto é fechado --- contém os seus
  limites.
\item \textbf{LEI.} Completude: diferença $\mathbb{Q}$ vs $\mathbb{R}$.
\item \textbf{TESTEMUNHA.} Convergentes de $\sqrt{2}$ em $[1,2]\cap\mathbb{Q}$:
  dentro, sem limite racional.
\item \textbf{CONCLUSÃO.} Em $\mathbb{R}$ o limite fica em $K$; em $\mathbb{Q}$ o
  mesmo intervalo falha.
\item \textbf{VOLTA.} Weierstrass herda: sem compacidade, contínua pode não atingir
  máximo.
\end{enumerate}
\medido{\code{analise\_resolve(14)}: $[1,2]\cap\mathbb{Q}$ não compacto.}

%----------------------------------------------------------------------
\section{heine borel}\label{sec:heine-borel}
\textbf{Enunciado.} Compacto $\Leftrightarrow$ fechado e limitado --- as duas
direcções.

\begin{enumerate}\itemsep2pt
\item \textbf{DEFINIÇÃO.} $K\subseteq\mathbb{R}$ (métrico).
\item \textbf{TRADUÇÃO.} $K$ compacto $\Leftrightarrow$ $K$ fechado e limitado.
\item \textbf{OPERAÇÃO.} «Compacto $\Rightarrow$ fechado e limitado» geral; a
  RECÍPROCA precisa de $\mathbb{R}^n$ e FALHA em $\mathbb{Q}$.
\item \textbf{LEI.} Completude sustenta a recíproca.
\item \textbf{TESTEMUNHA.} $[1,2]\cap\mathbb{Q}$: fechado em $\mathbb{Q}$, limitado,
  NÃO compacto.
\item \textbf{CONCLUSÃO.} «Fechado e limitado $\Rightarrow$ compacto» é de
  $\mathbb{R}^n$, não verdade geral.
\item \textbf{VOLTA.} Trio: uniforme (13), extremos (16), Heine--Borel (esta) --- uma
  hipótese.
\end{enumerate}
\medido{\code{analise\_resolve(15)}: contra-exemplo em $\mathbb{Q}$ fecha.}

%----------------------------------------------------------------------
\section{weierstrass}\label{sec:weierstrass-ar}
\textbf{Enunciado.} Contínua em compacto atinge máximo e mínimo.

\begin{enumerate}\itemsep2pt
\item \textbf{DEFINIÇÃO.} $f$ contínua no compacto $[a,b]$.
\item \textbf{TRADUÇÃO.} Contínua em $K$ compacto $\Rightarrow$ atinge máximo e
  mínimo.
\item \textbf{OPERAÇÃO.} «Atinge»: $f(K)$ limitado (há $\sup$) e fechado ($\sup$
  pertence) --- imagem contínua de compacto.
\item \textbf{LEI.} Heine--Borel na imagem.
\item \textbf{TESTEMUNHA.} $x^2-2x$ em $[0,3]$: extremos em pontos do compacto;
  $1/x$ em $(0,1)$: $\sup$ cresce com a malha.
\item \textbf{CONCLUSÃO.} Compacidade faz «atinge»; no aberto a contínua nem é
  limitada.
\item \textbf{VOLTA.} Trio do andar: uniforme / extremos / Heine--Borel --- três
  teoremas, uma hipótese.
\end{enumerate}
\medido{\code{analise\_resolve(16)} / \code{an\_extremos}: mín $-1$ em $x=1$ exacto.}

%----------------------------------------------------------------------
\section{derivada limite e quociente}\label{sec:derivada-ar}
\textbf{Enunciado.} O limite contra o quociente polinomial do Cálculo I.

\begin{enumerate}\itemsep2pt
\item \textbf{DEFINIÇÃO.} $f$ polinomial sobre $\mathbb{Q}$; $[a,b]$ racional.
\item \textbf{TRADUÇÃO.} $f'(a)=\lim_{h\to0}(f(a+h)-f(a))/h$ e $f'(a)=q(0)$.
\item \textbf{OPERAÇÃO.} Análise usa LIMITE; Cálculo I usa DIVISÃO EXACTA --- mesmo
  número; em polinómios $q(h)$ é polinómio.
\item \textbf{LEI.} Rolle (as outras derivam dela); $c$ pode ser irracional.
\item \textbf{TESTEMUNHA.} Três vias da derivada em muitos pontos; $0$ divergências.
\item \textbf{CONCLUSÃO.} As vias concordam; o $c$ do Valor Médio pode ser irracional
  (supremo em $\mathbb{R}$).
\item \textbf{VOLTA.} Cálculo I devolvia o corte; aqui há a maquinaria para o PROVAR.
\end{enumerate}
\medido{\code{analise\_resolve(17)}: três vias; $0$ divergências.}

%----------------------------------------------------------------------
\section{valor medio na analise}\label{sec:valor-medio-ar}
\textbf{Enunciado.} $f'(c)=(f(b)-f(a))/(b-a)$, e o $c$ pode ser irracional.

\begin{enumerate}\itemsep2pt
\item \textbf{DEFINIÇÃO.} $f$ polinomial; intervalo racional.
\item \textbf{TRADUÇÃO.} $f'(c)=(f(b)-f(a))/(b-a)$.
\item \textbf{OPERAÇÃO.} Rolle na função inclinada --- não é teorema novo.
\item \textbf{LEI.} Rolle; em polinómios tudo exacto excepto o ponto $c$.
\item \textbf{TESTEMUNHA.} $x^3-3x$ em $[0,2]$: $c=2/\sqrt{3}$ por encaixe.
\item \textbf{CONCLUSÃO.} $c$ pode ser IRRACIONAL --- vive como supremo.
\item \textbf{VOLTA.} Análise explica o que o Cálculo I já devolveu por bisseção.
\end{enumerate}
\medido{\code{analise\_resolve(18)}: $c$ irracional; encaixe.}

%----------------------------------------------------------------------
\section{rolle e o gume}\label{sec:rolle-ar}
\textbf{Enunciado.} Retirar $f(a)=f(b)$ e PROCURAR o contra-exemplo.

\begin{enumerate}\itemsep2pt
\item \textbf{DEFINIÇÃO.} Mesmo quadro polinomial.
\item \textbf{TRADUÇÃO.} $f(a)=f(b)\Rightarrow\exists c:\ f'(c)=0$.
\item \textbf{OPERAÇÃO.} Gume: sem $f(a)=f(b)$, o teorema CAI --- $f(x)=x$ em
  $[0,1]$, $f'\equiv1$.
\item \textbf{LEI.} A hipótese trabalha; o contra-exemplo é o mais barato.
\item \textbf{TESTEMUNHA.} Busca de $c$ com $f'(c)=0$ em $f(x)=x$: não acha.
\item \textbf{CONCLUSÃO.} Hipótese necessária; sem ela não há $c$.
\item \textbf{VOLTA.} Máquina de $\mathbb{R}$: $c$ como supremo, não só corte.
\end{enumerate}
\medido{\code{analise\_resolve(19)}: gume $f(x)=x$; busca vazia.}

%----------------------------------------------------------------------
\section{taylor}\label{sec:taylor-ar}
\textbf{Enunciado.} O resto medido, não estimado.

\begin{enumerate}\itemsep2pt
\item \textbf{DEFINIÇÃO.} $f=x^3-3x+2$, $a=1$.
\item \textbf{TRADUÇÃO.} $f(x)=\sum_{k=0}^{n}f^{(k)}(a)/k!\,(x-a)^k+R_n(x)$.
\item \textbf{OPERAÇÃO.} Medir o resto: construir $T_n$ e subtrair --- resto é
  função, não cota.
\item \textbf{LEI.} Coeficiente $k$ $=$ $f^{(k)}(a)/k!$ exacto em $\mathbb{Q}$.
\item \textbf{TESTEMUNHA.} $R_n$ anula-se a partir do grau; $R_n(a)=0$ em toda a
  ordem.
\item \textbf{CONCLUSÃO.} Em polinómios Taylor é IDENTIDADE --- resto $=$ parte cortada.
\item \textbf{VOLTA.} Fora dos polinómios o resto NÃO se anula: majorar é o conteúdo
  analítico.
\end{enumerate}
\medido{\code{analise\_resolve(20)} / \code{an\_taylor}: resto zero do grau em diante.}

%----------------------------------------------------------------------
\section{series}\label{sec:series-ar}
\textbf{Enunciado.} Critérios ESTRUTURAIS, e nunca pela soma parcial.

\begin{enumerate}\itemsep2pt
\item \textbf{DEFINIÇÃO.} $\sum1/n$, $\sum1/n^2$, $\sum(-1)^n/n$.
\item \textbf{TRADUÇÃO.} Decidir por critérios --- não pela soma parcial.
\item \textbf{OPERAÇÃO.} Blocos (harmónica); majoração telescópica ($p\ge2$);
  alternada.
\item \textbf{LEI.} Comparação com soma que TELESCOPA: $1/n^p\le1/(n-1)-1/n$.
\item \textbf{TESTEMUNHA.} Três veredictos sem somar; majoração falha exactamente em
  $p=1$; telescopagem $=$ forma fechada.
\item \textbf{CONCLUSÃO.} Critérios decidem sem monstro aritmético.
\item \textbf{VOLTA.} Forma fechada custa nada; soma bruta estoura --- medido.
\end{enumerate}
\medido{\code{analise\_resolve(21)}: blocos / majoração / telescopagem.}

%----------------------------------------------------------------------
\section{convergencia uniforme}\label{sec:conv-uniforme}
\textbf{Enunciado.} $x^n$ em $[0,1]$: pontual sim, uniforme NÃO.

\begin{enumerate}\itemsep2pt
\item \textbf{DEFINIÇÃO.} $f_n(x)=x^n$ em $[0,1]$.
\item \textbf{TRADUÇÃO.} $f_n\to f$ ponto a ponto, mas $\sup|f_n-f|\to0$ falha.
\item \textbf{OPERAÇÃO.} Mesma troca de quantificadores: $N$ por $x$ vs $N$ para
  todos; $N$ explode perto de $1$.
\item \textbf{LEI.} $\forall x\,\exists N$ contra $\exists N\,\forall x$.
\item \textbf{TESTEMUNHA.} $x_n=1-1/(2n)$ com $x_n^n>1/2$; limite pontual
  descontínuo em $1$.
\item \textbf{CONCLUSÃO.} Pontual sim, uniforme NÃO.
\item \textbf{VOLTA.} Uniforme $\Rightarrow$ pontual; a recíproca precisa de
  compacidade.
\end{enumerate}
\medido{\code{analise\_resolve(22)} / \code{an\_falha\_uniforme\_seq}: $x_n^n>1/2$.}

%----------------------------------------------------------------------
\section{riemann}\label{sec:riemann-ar}
\textbf{Enunciado.} $U(f,P)-L(f,P)\to0$, e a distância é exacta.

\begin{enumerate}\itemsep2pt
\item \textbf{DEFINIÇÃO.} $f=x^2$ em $[0,2]$; $P$ partição uniforme.
\item \textbf{TRADUÇÃO.} $U(f,P)-L(f,P)=(f(b)-f(a))\cdot h$ EXACTO.
\item \textbf{OPERAÇÃO.} Monótona: diferença TELESCOPA --- igualdade, não estimativa.
\item \textbf{LEI.} Telescopagem (mesma das séries).
\item \textbf{TESTEMUNHA.} $U-L$ coincide com $(f(b)-f(a))\cdot h$ em todas as
  partições; $\int$ no meio.
\item \textbf{CONCLUSÃO.} Integrabilidade pela igualdade --- $h\to0$ de graça.
\item \textbf{VOLTA.} Cálculo I: primitiva dava o valor; aqui o sanduíche explica POR
  QUÊ.
\end{enumerate}
\medido{\code{analise\_resolve(23)}: $U-L$ exacto; $0$ falhas.}

%----------------------------------------------------------------------
\section{troca de limite}\label{sec:troca-limite}
\textbf{Enunciado.} Qual hipótese torna $\lim\int=\int\lim$ legítimo.

\begin{enumerate}\itemsep2pt
\item \textbf{DEFINIÇÃO.} $f_n(x)=x^n$ em $[0,1]$.
\item \textbf{TRADUÇÃO.} $\lim\int f_n=\int\lim f_n$ sob QUE hipótese?
\item \textbf{OPERAÇÃO.} Hipótese suficiente: convergência UNIFORME --- aqui os lados
  coincidem SEM ela (cuidado: suficiente, não necessária).
\item \textbf{LEI.} $|\int f_n-\int f|\le\sup|f_n-f|\cdot(b-a)$ --- o SUP precisa de
  ir a zero.
\item \textbf{TESTEMUNHA.} $\int$ desce a $0$; $\sup$ fica em $1$; não-uniformidade
  exibida.
\item \textbf{CONCLUSÃO.} Uniforme é suficiente e NÃO necessária neste exemplo.
\item \textbf{VOLTA.} Hipótese mais fraca: dominação ($x^n\le1$) --- razão VERDADEIRA
  da troca aqui.
\end{enumerate}
\medido{\code{analise\_resolve(24)}: lados coincidem; $\sup$ não desce.}

%----------------------------------------------------------------------
\section{cinco vias}\label{sec:cinco-vias}
\textbf{Enunciado.} As cinco formas da completude --- e NENHUMA arbitra as outras.

\begin{enumerate}\itemsep2pt
\item \textbf{DEFINIÇÃO.} Cinco formas da completude de $\mathbb{R}$.
\item \textbf{TRADUÇÃO.} Dedekind $\cdot$ Cauchy $\cdot$ supremo $\cdot$ encaixantes
  $\cdot$ Bolzano--Weierstrass.
\item \textbf{OPERAÇÃO.} Cada uma SOZINHA, com testemunha da falha em $\mathbb{Q}$;
  NENHUMA função chama outra.
\item \textbf{LEI.} Equivalência demonstra-se --- não se escolhe árbitro.
\item \textbf{TESTEMUNHA.} Cinco colunas: passo de Möbius; módulo de Cauchy; Newton;
  bisseção; caixa --- mesmo buraco ($p^2\neq2q^2$).
\item \textbf{CONCLUSÃO.} Cinco falham em $\mathbb{Q}$ pelo MESMO motivo e por
  CAMINHOS INDEPENDENTES.
\item \textbf{VOLTA.} Por que $\mathbb{R}$ precisava de existir: cinco perguntas, a
  mesma resposta «não» em $\mathbb{Q}$; em $\mathbb{R}$ passam a «sim».
\end{enumerate}
\medido{\code{analise\_resolve(25)}: cinco programas; $0$ falhas; buraco comum.}

%======================================================================
\section{o que fica no conversa}
A prosa canónica AR1--AR25 é este ficheiro. \code{analise\_resolve} \textbf{mede}.
Catálogo de Análise Real fechado.

\end{document}
